% Generated by src/make_paper_tex.py from results/attack_table.csv.
At the same word budget, the random attack flips 28 of 1035 caught messages (2.7\%) and the saliency attack flips 46 of 1001 (4.6\%).
Each saliency-attacked message costs a mean of 23 model queries. The random word attack is not a single draw either: it re-queries every original and spends one query per attempt, so its 1{,}035 messages cost 3{,}927 queries, 3.8 each. The saliency attacker buys its extra flips with about 6.1 times as many calls, and that price is part of the claim, not a footnote to it.
The saliency attack also edits MORE, not less: 9.3\% of words against 7.6\% for the random attack. Normalising by the budget actually spent, it buys 0.493 of a flip per percent of words edited against 0.355, so the advantage is real but smaller than the raw ASR ratio.
Both numbers are pooled over different candidate vocabularies and, after the exclusions, over different message sets, so the denominator is 1{,}001 for the saliency attack and 1{,}035 for the random one; the paired comparison of Section~\ref{sec:pairs} is the only place where the two are measured on the same messages.
Pooled query accounting, kept separate from the message count because the two are different numbers: the character attack sends 1{,}044 caught messages in 4{,}120 queries, the word attack 1{,}035 in 3{,}927, so the random attacks together send 2{,}079 messages in 8{,}047 queries (3.9 each), and the saliency attack sends 1{,}001 messages in 22{,}964 queries (22.9 each).
